\subsection*{常用编译选项}

成功运行 \co@ele_add@ 后，就可以用同样的命令形式编译其他 \co@.co@ 文件。Croqtile 先生成目标源码，再调用目标工具链完成编译和链接。默认输出名为 \co@a.out@，使用 \co@-o@ 可以指定名称。

下表列出了选择输出形式和检查编译过程时常用的选项：

\begin{table}[htbp]\centering\small
\begin{tabularx}{\linewidth}{@{}lX@{}}
\toprule
\textbf{选项} & \textbf{作用} \\
\midrule
\texttt{-o \textless{}file\textgreater{}} & 指定输出文件名。 \\
\texttt{-t \textless{}platform\textgreater{}} & 选择目标平台；用 \texttt{--help-target} 查看支持的目标。 \\
\texttt{-es} & 只生成目标源码，不进行后续编译。 \\
\texttt{-E} & 只输出预处理后的 Croqtile 源码。 \\
\texttt{-c} & 编译为目标文件，不链接。 \\
\texttt{-S} & 生成汇编代码，需要目标工具链支持。 \\
\texttt{-gs} & 生成用于后续编译和执行的脚本。 \\
\texttt{-v} / \texttt{--verbose} & 显示编译器调用的外部命令。 \\
\texttt{--help} / \texttt{--help-hidden} & 查看常规或高级编译选项。 \\
\bottomrule
\end{tabularx}
\end{table}
